\section{Smooth degree-seven covers near a change of sign}\label{sec:global}
Changing the sum of the weights at a fixed set of flags through zero reverses the degree of the local smoothing equation. Smooth maps near the join account for this difference.
\subsection{Smooth covers approaching a joined map}\label{p14-sec:smooth}

Fix a simply branched cubic $Y\to C$ with monodromy $S_3$ and
$b=2\ell>0$ branch values. Let $I$ be a nonempty subset of these
values. Write $j=|I|$, $h_j=1+7(g-1)+3\ell+j$, and $f_I$ for the
finite map obtained by joining its Galois closure $Z$ to $C$ at
the selected flags fixed by one transposition. We use increasing
$t$ as the positive direction across a cubic wall. A crossing on a
smooth cover is counted once for each isomorphism class over $C$ and
each parameter, with the sign of the derivative of the normalized
kernel equation in the complex-oriented Hurwitz directions and the
oriented path direction.

\begin{theorem}\label{p14-thm:smooth-covers}
On the prescribed cubic and normal bundle there is a two-parameter
family of operators $D_{a,t}$ with the following properties.
Along $t=0$ the prescribed cubic has a normal kernel, transverse in
the $t$ direction. There is a sequence of parameters approaching a point
of this segment for which smooth connected degree-seven covers with a
nonzero normal section converge to $f_I$ and its root section. Their source genus is $h_j$ and their monodromy
is $S_7$. If $I$ consists of all branch values, the graphs in the
convergent sequence are embedded after discarding finitely many terms.

There are also arbitrarily small Hurwitz neighborhoods $U$ of $f_I$
and rectangles $[a_-,a_+]\times[-\delta,\delta]$ as above with the
following signed count. Let $d_L,d_R$ denote the local degrees in the
uniformizing chart of the joined map along the left and right edges.
On the lower and upper edges, oriented by increasing $a$, take
sufficiently small generic perturbations relative to their ends.
Their smooth crossings in $U$ are regular, have simple branching over
distinct values, trivial deck group and embedded kernel graph. Their
signed counts $N_-,N_+$ satisfy
\begin{equation}\label{p14-eq:smooth-edge}
                 2(N_+-N_-)=d_R-d_L\in\{-2,2\}.
\end{equation}
\end{theorem}

\begin{proof}
\emph{Changing the degree.}
Take a short segment $c(a)\in\mathbb R^b$ crossing the hyperplane
$\sum_{i\in I}c_i=0$ at a point of no other subset hyperplane.
In the branch coordinates of Section~\ref{p14-sec:joins}, $H_i$ is
a nonzero real functional on a complex smoothing plane. Choose
nonzero $\gamma_i(a)$ smoothly with $H_i(\gamma_i(a))=c_i(a)$.
A right inverse of $H_i$ plus a fixed nonzero vector in its kernel
gives such a choice even when $c_i=0$. Assign the antilinear
coefficient at the $i$th center so that
$2\overline{\gamma_i(a)}[B_iv_i]=[e_i]$.
The local replacement and restoration on a disk of regular values in
Section~\ref{p14-sec:joins} apply with one cutoff radius for this compact
segment. Indeed the prescribed coefficients and the right inverses
in the restoration are uniformly bounded. Thus the weights on the
restored cubic are uniformly close to $c_i(a)$. Take the radius
small enough to retain opposite signs of the $I$-sum at the two ends
and to keep every other subsum away from zero. The ordinary and polar
operators on the fixed covers retain the ranks used in the nodal
derivative. The cubic remains transverse to one common time direction.
The kernel at $f_I$ remains the odd real line: evaluation at any
chosen flag removes the even section on $Z$.

Write $q_i$ for the complex smoothing parameters, and $h$ for the
remaining chart coordinates and time, centered at $f_I$, $t=0$.
Trivialize the reduced obstruction equation over the parameter
segment by stabilizing on disks of regular values and normalizing the
root section by a real linear functional. The center kernel and
cokernel have constant rank, so supplements and normalizations can be
continued along the segment. By Proposition~\ref{p14-prop:signs} the end
zeros are isolated and have opposite degrees in this uniformizing
chart. On the boundary of every sufficiently small ball in $(h,q)$
a zero occurs at some intermediate $a$, by homotopy invariance of
degree. Passing to a sequence of shrinking balls and a convergent
subsequence of parameters, put $r=\max_{i\in I}|q_i|$. The nodal
derivative is uniformly injective, hence $r>0$ at every noncentral
zero sufficiently close to the center and $|h|=O(r\log(1/r))$.
After division by $r\log(1/r)$ the matching equation of
Section~\ref{p14-sec:joins} places every subsequential limit of $(q_i/r)$
on the real line spanned by $(\beta_i)_{i\in I}$ at the limiting
parameter. Here $2\overline{\beta_i}[B_i v_i]=[e_i]$. Every coordinate
of the generator is nonzero. Thus all $|q_i|$ are comparable to $r$.
The remaining adjoint row, divided by $r$, now gives
$\sum_{i\in I}H_i(\beta_i)=0$ at the limit. These estimates are
uniform on the segment, since they use the bounds for the stabilized
inverse and the leading nodal and matching maps, all of which retain
their ranks at resonance. Every selected node has smoothed. The
domain is connected, of arithmetic genus $h_j$, and the normalized
sections converge to the odd section. This proves existence of the
asserted sequence.

\emph{Ramification and separation.}
On the three local sheets involved at a selected flag, the projection
in the neck coordinate $z$ is
\[
              w-c_i'=(z+p_i)^2+q_i/z .
\]
Here $c_i'$ is the displacement of the branch center.
Both $c_i'$ and $p_i$ are $O(r\log(1/r))$. Put $d_i^3=q_i/2$.
The critical equation is $2z^2(z+p_i)=q_i$, with three distinct
simple roots $z/d_i=\zeta+o(1)$, $\zeta^3=1$. At such a root,
\[
        w-c_i'=(z+p_i)(3z+p_i)=3d_i^2(\zeta^2+o(1)).
\]
These critical values have mutual distances comparable to $r^{2/3}$.
The two other local quadratic factors have critical values at
distance $O(r\log(1/r))$ from the original branch value. Consequently
each of the three displayed values supports exactly one simple
ramification point. At each unselected branch disk the three
quadratic factors persist. In total the ramification degree is
$B=3b+2j$. A solitary simple branch gives a transposition on the
seven sheets. The graph of the conjugates of that transposition has
components permuted transitively by the monodromy, so their common
size divides seven. There is an edge, hence this graph is connected
and its edge transpositions generate $S_7$.

For the full join the two limiting images coincide only at the
selected points. On a selected disk the normalized height on the
cubic factor has the expression $U_0(w)+zU_1(w)+z^2U_2(w)$. In the
fixed-base polynomial presentation its coefficients converge in
$W^{1,p}$ for some $2<p<4$, hence uniformly, by the stabilized
inverse in Section~\ref{p14-sec:joins}. Centrally $U_1(0)=v_i\ne0$ and
$U_2=0$. The difference of the heights at distinct roots $z,z'$
over the same $w$, divided by $z-z'$, is
\[
                          U_1(w)+(z+z')U_2(w).
\]
Its norm is bounded below on a smaller disk. At a simple critical
point with local coordinate $w=w_0+\eta^2$, the smooth height has
zero antiholomorphic derivative at $\eta=0$, by the pullback
equation. Taking the difference at $\eta$ and $-\eta$ gives vertical
derivative in the $z$ coordinate equal to
$U_1(w_0)+2z_0U_2(w_0)\ne0$. Thus this factor also immerses at
ramification. The other two quadratic factors stay separated by
their distinct nonzero limiting heights. Their normalized heights
converge in $C^1$ on smaller source disks by ordinary elliptic
regularity and have nonzero limiting vertical derivative at the
ramification points. Away from these disks the projections are
unramified and the central heights are pairwise separated on a
compact set. The resulting injective immersion of a compact source is
an embedding.

\emph{Isolating a neighborhood for paths of operators.}
To obtain signed counts, the neighborhood must lie in Hurwitz space
alone, independent of $a$ and $t$. Arrange the infinite-order Petri
condition at two distinct points of a regular disk where two cubic
heights form a complex frame, by the local gauge construction of
Section~\ref{p14-sec:kernels} and \cite[Corollary~5.10]{wendl}. Restore
the cubic, and choose its transverse time direction, on another
subdisk. Both Petri jets are then fixed along the family. Vary the
central segment relative to its ends on the frame disk by small
$a$-dependent tensors flat at the test points and annihilating the two
current cubic heights. This preserves the prescribed cubic.

For connected smooth covers sufficiently close to $f_I$, monodromy
outside the branch disks contains the regular $S_3$ on six sheets
fixing the seventh. Connectedness now implies double transitivity.
The reduced real sheet representation $W$ is absolutely irreducible
of dimension six: the complex commuting algebra of the permutation
representation has dimension two by the two orbits on ordered pairs.
The constant sector is invertible. The six components of a nonzero
solution in $W$ and the two independent cubic heights are linearly
independent as real solution germs on a regular disk. Indeed on the
joint normal cover they span inequivalent irreducible systems, and
unique continuation identifies their linear relations on the disk.

Here is the constrained transversality argument. An adjoint
annihilating the allowed variations gives on the disk a tensor
identity
\[
           \sum_{k=1}^{6}u_k\otimes v_k
                    =s_1\otimes\alpha_1+s_2\otimes\alpha_2,
\]
where $u_k,v_k$ are the test solution and adjoint components and the
$s_i$ are the prescribed cubic heights. The smooth frame gives the
$\alpha_i$ smoothly across the test points. Pair with $[D,\Phi]$ for
compactly supported complex-linear endomorphisms $\Phi$.
The left integral vanishes by the solution and adjoint equations.
Since $\Phi s_1,\Phi s_2$ can be chosen freely, both $\alpha_i$ solve
the adjoint equation. The Petri condition on infinite jets and
independence of the eight solution germs force the jets of all
$v_k$ to vanish. Unique continuation gives $v_k=0$. This is the
local-system Petri argument of \cite[\S\S1.6, 2.7--2.9]{dwbn}
with the two prescribed solutions retained.

A smooth branching profile with $s$ distinct values has complex
dimension $s$. The pullback operator has real index $-2B$ by
Riemann--Roch. After allowing $a$ and dividing by real kernel
scalings the expected dimension is $2s-2B$. The preceding test
gives transversality for a generic constrained central segment.
To apply it on the noncompact smooth strata, exhaust their
projective incidences by countably many compact subsets. Finitely
many supported variations suffice near the zeros on each such set,
and their parameter cutoffs preserve the end collars. The regularity
conditions on compact subsets are open, so the countable
intersection gives regular segments. Only simple branching at
distinct values remains, in dimension zero. The uniform matching
estimate above excludes partially smoothed zeros sufficiently near
the center throughout the segment. Thus all off-center zeros at
$t=0$ sufficiently near $f_I$ form a countable set.

Choose arbitrarily small invariant Hurwitz balls $U$ with radii
outside the set of radii of these zeros. Compactness gives a
$\delta>0$ for which $\partial U$ is free of kernels throughout the
rectangle. The uniform estimates exclude all nonsmooth zeros on
the horizontal faces after shrinking the ball and $\delta$, since
a zero with $q=0$ can only be the center at $t=0$.
The sign condition at each end leaves only the joined zero on each
vertical face. Small perturbations of the horizontal paths preserve
these boundary properties. Applying the preceding transversality
test there with unrestricted operator variations gives regular
zeros, with distinct simple branch values and one-dimensional
kernel. There is a transposition, hence full monodromy by the prime
degree argument. The deck group, the centralizer in $S_7$, is trivial.

For completeness, the marked tests for embeddedness are also
transverse. At two distinct preimages over a common value, equality
of heights has real codimension four, while a free base value has
two real parameters and a branch value is already determined.
Vanishing of the vertical complex derivative at a ramification
point gives another four real conditions. For either test a
putative annihilator solves its adjoint equation away from the
marked fiber. Its invariant part vanishes distributionally by base
invertibility: each test vanishes on sections pulled back from the
base. One of the two Petri points lies away from the marked value.
The unrestricted adjoint test there and unique continuation kill the
reduced part away from the marks. To test the remaining point-supported
distributions, vary a local section with prescribed value or
holomorphic first derivative. Its equation can be made zero to
any given finite order, since the coefficients involving a
$\bar\partial$ derivative are determined recursively by its pure
holomorphic jet. These variations kill the marked multipliers,
and the remaining homogeneous adjoint vanishes by elliptic regularity.
At a zero of a marked test it vanishes also on the infinitesimal
rescaling. The resulting negative expected dimensions exclude both
coincidences and nonimmersions. The signed zero sets on the compact
horizontal faces are therefore finite and have embedded graphs.

\emph{Boundary orientations.}
Orient the parameter rectangle by $da\wedge dt$. On the product of
its boundary with the uniformizing Hurwitz ball, the total signed
degree is zero because the section extends over the product and has
no zeros on the Hurwitz boundary. The right and left faces contribute
$d_R$ and $-d_L$. On each horizontal face a smooth zero has two
lifts, interchanged by the central involution. Their signs agree.
Indeed this involution acts complex linearly on the Hurwitz chart.
On the normalized obstruction target it preserves orientation,
as follows by equivariance from the nonzero local degree at either
vertical end. This orientation character persists across the
rectangle. The lower and upper faces thus contribute $2N_-$ and
$-2N_+$, respectively. This proves \eqref{p14-eq:smooth-edge}.
\end{proof}

When $I$ contains every branch value, complementary nonempty proper
subsets have opposite sums at the limiting weights. In particular the
size-$b/2$ terms cancel in either adjacent chamber. Weights of one
sign give the largest sum of absolute changes over all genera and
belong to different chambers. Equation~\eqref{p14-eq:smooth-edge}
determines the difference of the two counts in $U$. Extending that
identity to counts in the whole Hurwitz space would require control
of its complement.
